Our strategy for proving this theorem is to prove facts about the quivers $Q(\mathcal{T})$ for triangulations $\mathcal{T}$ of $C(n + 2d + 1, 2d)$ without internal $(d + 1)$-simplices. We then use these properties to show that one can realise $Q(\mathcal{T})$ as a slice of $\tqdn$, which will imply that $Q(\mathcal{T})$ is a cut of $\qdn$. The idea of the proofs of both of the following lemmas is that if $Q(\mathcal{T})$ does not possess a certain property, then we can find an interior $(d + 1)$-simplex in~$\mathcal{T}$.

\begin{lemma}\label{lem:single-step-arrows}
If $\mathcal{T}$ has no interior $(d+1)$-simplices, then the arrows in $Q(\mathcal{T})$ are all of the form $A \rightarrow A + E_{i}$.
\end{lemma}
\begin{proof}
By Proposition~\ref{prop:all_but_one_vertex}, every arrow is of the form $A \rightarrow A + \rvec$ for some $r>0$ and for each such arrow, we have that $(a_{0}, a_{1}, \dots, a_{i}, a_{i} + r, a_{i+1}, a_{i + 2}, \dots, a_{d})$ is a face of a $2d$-simplex of $\mathcal{T}$. If $r>1$, then this is an interior $(d+1)$-simplex.
\end{proof}

\begin{lemma}\label{lem:simp_arrow_switch}
Suppose that $\mathcal{T}$ has no interior $(d+1)$-simplices. If $A \rightarrow A + E_{i} \rightarrow A + E_{i} + E_{j}$ is a sequence of arrows in $Q(\mathcal{T})$, then, if $A + E_{j} \in \nonconsec{n+2d+1}{d}$, the sequence $A \rightarrow A + E_{j} \rightarrow A + E_{i} + E_{j}$ is also in $Q(\mathcal{T})$.
\end{lemma}
\begin{proof}
We assume $A + E_{j} \in \nonconsec{n+2d+1}{d}$. By re-ordering, we may assume that $i = d$. Then there is a $2d$-simplex $T$ of $\mathcal{T}$ such that $e(T) = A + E_{j} + E_{d}$ by Lemma~\ref{lem:find_2d_simplex}.

Let $T = (a_{0}, x_{1}, a_{1}, \dots, x_{j}, a_{j} + 1, x_{j + 1}, \dots, x_{d}, a_{d} + 1)$. If $x_{d} = a_{d}$, then $A + E_{j}$ is a $d$-face of $T$ and hence a $d$-arc of $\mathcal{T}$. Hence, suppose for contradiction that $a_{d} \neq x_{d}$, so that $a_{d - 1} < x_{d} \leqslant a_{d} - 1$. Note that if $x_{j} = a_{j - 1} + 1$, then $(a_{0}, a_{1}, \dots , a_{d}) \swr (x_{1}, x_{2}, \dots ,x_{d}, a_{d} + 1)$, which is a $d$-face of $T$. Hence $a_{j - 1} + 2 \leqslant x_{j} < a_{j} + 1$. Then $(a_{0}, a_{1}, \dots, a_{j - 1}, x_{j}, x_{j + 1}, \dots , x_{d}, a_{d} + 1)$ is a $(d + 1)$-face of $T$ and an interior $(d+1)$-simplex of $\mathcal{T}$, a contradiction.
\end{proof}
